\noindent The $T=60$ case is cleanest:
$\delta \in \{0, 2\pi/5, 4\pi/5\}$, with denominator exactly $L=5$.
This suggests a Diophantine condition $L_x\,\delta_x/\pi \in
\mathbb{Q}$ (and likewise for $y$) governing which phases permit
revival.

Dukes' 1D analysis established revival conditions on a $k$-cycle in
terms of $\cos(2\pi/k)$ and related algebraic quantities. Three
differences stand out in the 2D separable case:

\begin{enumerate}[nosep]
\item \textbf{Separability does not trivialise the problem.} The
revival condition couples $L_x$ and $L_y$ non-trivially; it is not a
product of two independent 1D conditions.

\item \textbf{The 2D period is not $\mathrm{lcm}(T_x,T_y)$.} For
example, $3\times4$ with $\rho_x=2/3$, $\rho_y=1/2$ revives at $T=8$,
which is not the lcm of any obvious pair of 1D periods.

\item \textbf{New algebraic values appear.} Values such as
$(2-\sqrt2)/4 \approx 0.146447$ do not arise as $\cos(2\pi/k)$ for
small $k$, indicating genuinely 2D algebraic structure.
\end{enumerate}
